\subsection{Removing the rational character of order six}
\label{sec:kt-order-six}

The class-number bounds used here are known results.  This section
records their consequence for the precise rational class component in
the report.  It uses the general relative-unit bound of Gras and the
explicit even-character bound of Louboutin; it does not import an
unread statement from the title of Moser's 1981 paper.

\begin{proposition}[KT-ORDER-SIX-COMPONENT-REDUCTION]
Let $p>3$ be prime and let $e\ge2$ be even.  If
$\operatorname{ord}(\omega^e)=6$, then
\[
                 A_p^{[-e]}=0,
                 \qquad K_{2e}(\mathbf Z)\{p\}=0.       \tag{OS6.1}
\]
There is no bound on $e$ relative to $p$.
Together with KT-SMALL-ORDER-COMPONENT-REDUCTION, a remaining
rational character has order $5$ or at least $7$.
\end{proposition}

\begin{proof}
Put $L=\mathbf Q(\mu_p)$ and
$F=L^{\ker(\omega^{-e})}$.  Then $F$ is the real cyclic
sextic field of conductor $p$; in particular $p\equiv1\pmod{12}$.
Write $F_2,F_3$ for its quadratic and cubic subfields,
and $h_i,R_i$ for their ordinary class numbers and regulators,
with $h_6=h(F)$ and $R_6=R(F)$.  All three conductors are $p$,
and the respective discriminants are $p,p^2,p^5$.

For $p<2^{31}$, the carried Hart--Harvey--Ong verification
gives $A_p^+=0$, so the assertion follows from the evenness
of $e$ and K17-(K2).  Hence assume $p\ge2^{31}$.

First,
\[
                         J:=\frac{h_6}{h_2h_3}
                            \quad\hbox{is a positive integer}. \tag{OS6.2}
\]
This is KT-SEXTIC-CLASS-PRODUCT, whose proof uses the ordinary
Hilbert class fields and total ramification at $p$.
The same divisibility is stated in Louboutin, printed p.~455,
immediately after Lemma~3.  The integrality in (OS6.2),
not just a real-valued class-number quotient, is used below.

Let $\sigma$ generate $\operatorname{Gal}(F/\mathbf Q)$,
and use the relative-unit conventions
\[
\begin{split}
 U^*&=\{u\in\mathcal O_F^\times:
       N_{F/F_2}(u),N_{F/F_3}(u)\in\{\pm1\}\},\\
 R^*&=(\log|\varepsilon_*|)^2+
       (\log|\varepsilon_*^\sigma|)^2
       -(\log|\varepsilon_*|)(\log|\varepsilon_*^\sigma|),
\end{split}                                                   \tag{OS6.3}
\]
where $\{-1,\varepsilon_*,\varepsilon_*^\sigma\}$ generates
$U^*$.  Louboutin, Lemma~3, gives
\[
                 R_6=\frac{12R_2R_3R^*}{Q},
                 \qquad Q\in\{1,3,4,12\}.             \tag{OS6.4}
\]
There is no extra signature index in $R^*$.
Indeed $N_{F/F_3}(u)=-1$ would give total norm $-1$,
whereas the norm through $F_2$ of $\pm1$ is $+1$.
Thus $N_{F/F_3}(u)=+1$, and replacing $u$ by $-u$
if necessary makes $N_{F/F_2}(u)=+1$ as well.
This identifies $U^*/\{\pm1\}$ with the exact norm-one
module $E_*$ in Gras, without changing absolute logarithms.

Gras, \S2(b), Theorem~(2.5), with its definition~(2.4),
therefore gives
\[
                    R^*\ge\frac3{16}
                    \log^2\!\left(\frac{p-20}{4}\right).
                                                               \tag{OS6.5}
\]
The bound applies to every real cyclic sextic field of
conductor at least $28$, not just to the special family
introduced later in her \S3.  Its normalization can also
be checked directly from her Lemma~(2.6): for the logarithms
$x,y,y-x,-x,-y,x-y$ of a relative unit, put
$m=\max(|x|,|y|,|x-y|)$.  The lemma gives
$2m\ge\log((p-20)/4)$, and
$x^2+y^2-xy\ge3m^2/4$.  Apply this to a generating unit
to obtain (OS6.5).

Let $\chi$ be either primitive even Dirichlet character of
order six associated with $F$.  The analytic class-number
formula and the factorization
$\zeta_F\zeta_{\mathbf Q}/(\zeta_{F_2}\zeta_{F_3})
=L(s,\chi)L(s,\overline\chi)$ give
\[
                  J=\frac{Qp}{48R^*}|L(1,\chi)|^2.
                                                               \tag{OS6.6}
\]
For example, the residues of the three zeta functions are
$2h_2R_2/\sqrt p$, $4h_3R_3/p$ and
$32h_6R_6/p^{5/2}$; substitution of (OS6.4) yields exactly
(OS6.6).  This also agrees with Louboutin's displayed
class-number identity on printed p.~456 after division
by the cubic class-number formula.

Louboutin, \emph{Acta Arith.}~101 (2002), Theorem~1(1),
equation~(3), gives, for every primitive even character,
\[
 |L(1,\chi)|\le\tfrac12(\log p+\kappa_0),
 \qquad\kappa_0=2+\gamma-\log(4\pi)<1.                \tag{OS6.7}
\]
Combining (OS6.5)--(OS6.7) and $Q\le12$ gives
\[
 J\le\frac p3
       \left(\frac{\log p+\kappa_0}
       {\log((p-20)/4)}\right)^2
   <\frac{49}{108}p<p.                                \tag{OS6.8}
\]
Here the last comparison has ample exact margins: set
$t=\log p>20$, using $\log2>2/3$.  Since $p>100$,
$(p-20)/4>p/5$, and $\log5<2$; hence the displayed
ratio is less than $(t+1)/(t-2)<21/18$.
The bound $\kappa_0<1$ follows from $\gamma<1$ and
$\log(4\pi)>2$.
This proves (OS6.8) without relying on a decimal
evaluation of Euler's constant.

The already imported quadratic and cubic bounds give
$h_2\le\sqrt p/2<p$ and $h_3<2p/3<p$.
Together with (OS6.2) and (OS6.8), they imply
$p\nmid h_6=h_2h_3J$.

Finally put $H=\operatorname{Gal}(L/F)$.
For extension and norm on ideal classes,
\[
                       jN(x)=\sum_{g\in H}g x.         \tag{OS6.9}
\]
If $x\in A_p^H$, then $N(x)=0$ because
$\operatorname{Cl}(F)\{p\}=0$.  Thus $|H|x=0$;
as $p\nmid|H|$, one has $A_p^H=0$.
The $\omega^{-e}$ component lies in these invariants,
so $A_p^{[-e]}=0$.  The carried comparison K17-(K2)
gives the stated $K$-group vanishing in every weight
with this periodic character.
\end{proof}

\paragraph{Exact primary inputs and read scope.}
M.-N.~Gras, \emph{Familles d'unit\'es dans les extensions
cycliques r\'eelles de degr\'e 6 de $\mathbf Q$},
Besan\c con 1984/85--1985/86, fascicule~2:
\S2(a), printed pp.~2--3, fixes the general-field and unit
conventions; \S2(b), definition~(2.4), Theorem~(2.5)
and Lemma~(2.6), printed p.~4, give the regulator bound,
with their complete proofs on printed pp.~4--6.
These are physical PDF pages~5--7 including the cover.
The general result precedes the special-family \S3.
The primary scan is
\url{https://www.numdam.org/article/PMB_1986___2_A2_0.pdf}.
The cited earlier resolvent-divisibility and general
lattice results remain named published inputs; no fresh
review of all of those earlier foundational proofs is claimed.

S.~Louboutin, \emph{Class numbers of real cyclotomic fields},
\emph{Publ. Math. Debrecen}~64 (2004), 451--461,
\S3.1, printed pp.~454--456, gives the relative-unit
definition, Lemma~3 and its proof, the class-number
divisibility, and the exact analytic identity used here:
\url{https://publi.math.unideb.hu/paper/1169/download/10_5486_PMD_2004_2976.pdf}.
His separate analytic lower bound in Theorem~5 is not used.

S.~Louboutin, \emph{Explicit upper bounds for $|L(1,\chi)|$
for primitive even Dirichlet characters},
\emph{Acta Arith.}~101 (2002), 1--18:
Theorem~1(1), equation~(3), printed p.~2, is (OS6.7).
The full \S3 proof, printed pp.~5--8, including Lemma~7
and the $r=0$ completion, was read; its hypotheses allow
nonreal characters of order six.  The official PDF is
\url{https://www.impan.pl/shop/publication/transaction/download/product/82711}.
Exact downloaded-source hashes and review scopes are
recorded in the accompanying evidence manifest.

The small-prime Hart--Harvey--Ong verification, the lower-degree
class-number bounds, and K17-(K2) remain the already accepted
inputs.  The degree-six class-number conclusion is known in the
literature; the contribution here is its explicit source-checked
application to the remaining rational character, not a claim to
a new degree-six class-number theorem.
